\documentclass{article}
\usepackage{graphicx} % Required for inserting images

\title{TPC8v2}
\author{Afonso V. Francisco }
\date{3 Maio 2024}
\usepackage{amsfonts}
\begin{document}

\maketitle

\textbf{Seja $\mathbb{R}$ o anel do números reais e seja $\mathbb{C}$ o anel dos números complexos. Determine geradores para o núcleo do homomorfismo $T: \mathbb{R}[x] \rightarrow \mathbb{C}$ dado por $T(f(x))=f(2+i)$.}\\

Temos $ker \ T= \{f(x) \in \mathbb{R}[x] \mid f(2+i)=0\}$.

Mas $f(2+i)=0$ significa que $2+i$ é uma raíz de $f$, logo $x-(2+i) \mid f(x)$. Mas, como $2+i$ é uma raíz o seu conjugado, $2-i$, também é uma raíz de $f$, logo $x-(2-i) \mid f(x)$.

Assim $(x-2-i)(x-2+i)=x^2-4x+5 \in \mathbb{R}[x]$ divide  $f(x)$. Ora, isso quer dizer que existe $q(x) \in \mathbb{R}[x]$ tal que $f(x)=(x^2-4x+5)q(x)$. Portanto $$ker\ T=(x^2-4x+5)$$




\end{document}